\subsection{Признаки абсолютной сходимости рядов (Коши и Даламбера)}

\subsubsection*{Теорема. Признак Коши (радикальный)}
Пусть дан ряд $\sum_{n=1}^\infty a_n$. Вычислим верхний предел:
\[ q = \overline{\lim}_{n\to\infty} \sqrt[n]{|a_n|} \]
Тогда:
\begin{enumerate}
    \item $q < 1 \Rightarrow$ ряд сходится абсолютно.
    \item $q > 1 \Rightarrow$ ряд расходится.
    \item $q = 1 \Rightarrow$ неопределенность (может сходиться или расходиться).
\end{enumerate}

\begin{tcolorbox}[title=Доказательство]
    \textbf{1. Сходимость ($q < 1$):}
    Выберем число $p$ так, что $q < p < 1$.
    По определению верхнего предела, начиная с некоторого номера $N$, все члены последовательности $\sqrt[n]{|a_n|}$ будут меньше $p$:
    \[ \exists N: \forall n > N \quad \sqrt[n]{|a_n|} < p \iff |a_n| < p^n \]
    Ряд $\sum p^n$ — сходящаяся геометрическая прогрессия ($p < 1$).
    По признаку сравнения, $\sum |a_n|$ сходится $\Rightarrow \sum a_n$ сходится абсолютно.

    \textbf{2. Расходимость ($q > 1$):}
    Если верхний предел $q > 1$, то существует подпоследовательность $n_k$, для которой:
    \[ \sqrt[n_k]{|a_{n_k}|} > 1 \iff |a_{n_k}| > 1 \]
    Значит, общий член ряда не стремится к нулю ($a_n \not\to 0$).
    Нарушено необходимое условие сходимости $\Rightarrow$ ряд расходится.
\end{tcolorbox}

\textbf{Замечание.} Признак Коши является \textbf{достаточным} условием. Он эффективен, когда общий член содержит степень $n$.

\textbf{Пример 1.} $\sum_{n=1}^\infty (2+(-1)^n)^n \cdot z^n$
\[ \sqrt[n]{|a_n|} = (2+(-1)^n)|z| \]
\begin{itemize}
    \item При четном $n=2k$: $\sqrt[n]{|a_n|} = 3|z|$.
    \item При нечетном $n=2k+1$: $\sqrt[n]{|a_n|} = |z|$.
\end{itemize}
Верхний предел $q = \overline{\lim} \sqrt[n]{|a_n|} = 3|z|$.
\begin{itemize}
    \item Если $3|z| < 1 \Rightarrow |z| < 1/3$ — сходится.
    \item Если $3|z| > 1 \Rightarrow |z| > 1/3$ — расходится.
\end{itemize}

\subsubsection*{Теорема. Признак Даламбера}
Пусть существует предел (конечный или бесконечный):
\[ q = \lim_{n\to\infty} \left| \frac{a_{n+1}}{a_n} \right| \]
Тогда:
\begin{enumerate}
    \item $q < 1 \Rightarrow$ ряд сходится абсолютно.
    \item $q > 1 \Rightarrow$ ряд расходится.
    \item $q = 1 \Rightarrow$ неопределенность.
\end{enumerate}

\begin{tcolorbox}[title=Доказательство]
    \textbf{1. Сходимость ($q < 1$):}
    Выберем $p$ так, что $q < p < 1$.
    Начиная с некоторого номера $N$:
    \[ \forall n \ge N: \quad \left| \frac{a_{n+1}}{a_n} \right| < p \implies |a_{n+1}| < p |a_n| \]
    Распишем для произвольного $n > N$:
    \[ |a_n| < p |a_{n-1}| < p^2 |a_{n-2}| < \dots < p^{n-N} |a_N| \]
    \[ |a_n| < \frac{|a_N|}{p^N} \cdot p^n = C \cdot p^n \]
    Ряд мажорируется геометрической прогрессией $\sum C p^n$, значит, сходится.

    \textbf{2. Расходимость ($q > 1$):}
    Начиная с некоторого номера $N$:
    \[ \left| \frac{a_{n+1}}{a_n} \right| > 1 \implies |a_{n+1}| > |a_n| \]
    Последовательность модулей возрастает, значит $a_n \not\to 0$. Ряд расходится.
\end{tcolorbox}

\textbf{Важный факт:} Признак Коши \textbf{сильнее} признака Даламбера.
Если предел в признаке Даламбера существует, то признак Коши даст тот же результат. Но бывают случаи, когда Даламбер не работает (предела нет), а Коши работает.

\textbf{Пример 2 (Коши работает, Даламбер нет).}
Ряд: $\sum_{n=1}^{\infty} \frac{1}{2^n}$ с переставленными степенями множителей (для наглядности возьмем пример из вашего конспекта):
\[ a_n = \frac{1}{(3+(-1)^n)^n} \]
\begin{itemize}
    \item \textbf{Коши:}
    $\sqrt[n]{|a_n|} = \frac{1}{3+(-1)^n}$. Значения чередуются: $1/4, 1/2, 1/4, 1/2 \dots$
    Верхний предел $q = 1/2 < 1$. \textbf{Ряд сходится.}
    
    \item \textbf{Даламбер:}
    Рассмотрим отношение $D_n = \left| \frac{a_{n+1}}{a_n} \right|$.
    \begin{itemize}
        \item Если $n=2k$: $\frac{a_{2k+1}}{a_{2k}} = \frac{(1/2)^{2k+1}}{(1/4)^{2k}} = \frac{4^{2k}}{2 \cdot 2^{2k}} = \frac{2^{4k}}{2^{2k+1}} = 2^{2k-1} \to \infty$.
        \item Если $n=2k+1$: $\frac{a_{2k+2}}{a_{2k+1}} = \frac{(1/4)^{2k+2}}{(1/2)^{2k+1}} = \frac{2^{2k+1}}{4^{2k+2}} \to 0$.
    \end{itemize}
    Предел отношения не существует (скачет от 0 к $\infty$). Признак Даламбера неприменим.
\end{itemize}

